\section{结论}

\begin{enumerate}[leftmargin=2em,itemsep=0.35em]
  \item 建立了变截面矩形拉杆的三种自写有限元方法，并用 CalculiX CPS4/CPS8 与 FEniCS 完成第三方锚点验证。同网格 Q4 与 CPS4 的端位移差不超过 $0.05\%$，FEniCS 对 S1 的端位移与 CCX 参考仅差 $+0.0029\%$，证明单元、组装、边界与载荷实现正确。
  \item 一维杆采用柔度积分等效刚度后，节点位移与解析解一致到 $10^{-13}$ 且与网格无关；平均面积刚度会高估刚度，粗网格端位移误差按 $O(h^2)$ 收敛。等截面退化校核（$\beta=0$）通过。
  \item Q4（$2\times2$ 全积分）在 $160\times32$ 下的位移偏差不超过 $0.025\%$、探针应力偏差不超过 $0.042\%$；位移收敛阶约 $1.7$–$2.5$，应力收敛阶约 $1.6$–$2.3$，梯度保真误差 $0.35\%$–$0.95\%$。
  \item CST 的全局量在对称工况下与剖分方向无关（离散度约 $10^{-11}\%$），局部单元应力对角离散度为 $0.2\%$–$1.1\%$；面积加权节点恢复在粗网格上显著降低截面内应力振荡，但不能替代网格加密。
  \item 一维判据 $\eta=|\beta|h_0/L$ 与二维参考一致：S0/S1/S2 为 A，S3 为 B，S4 为 C；一维杆端部偏差依次为 $+0.25\%$、$-0.55\%$、$+0.05\%$、$-3.21\%$、$-7.36\%$，中截面偏差为 $+0.50\%$、$-0.05\%$、$+0.29\%$、$-1.16\%$、$-4.10\%$。
  \item 推荐流程：A/B 级工况先用等效刚度杆单元估算轴力与位移，再用 Q4 复核端部；C 级工况必须使用二维模型；CST 适合快速原型/复杂几何，但需注意局部应力对剖分的敏感性。三种方法的产出形态差异（一维曲线、连续云图、分块云图）应作为选型的首要判据之一。
\end{enumerate}
